\bookchapter{제21장: 모바일 네트워크와 정보 보안 심층 탐구}

\chapterintro{우리는 매일 스마트폰으로 인터넷, 결제, 대화를 하며 매우 편리합니다. 하지만 생각해 본 적 있나요? 채팅 기록, 결제 비밀번호, 사진... 이런 정보가 인터넷에서 “벌거벗고” 노출될 일이 얼마나 될까요? 솔직히 대부분의 사람들은 이런 문제를 생각해 본 적도 없지만, 문제가 터져서 — 예를 들어 계정이 도난당하거나 은행 카드가 사용된다면 — 그때 후회해도 소용없습니다. 이 장에서는 디지털 라이프를 덜 “투명하게” 만드는 방법을 자세히 이야기해 봅니다.}

\booksection{하나, VPN: 네트워크의 “투명 망토”}

VPN은 Virtual Private Network(가상 사설망)의 약자입니다. 여러분의 집과 가고 싶은 웹사이트 사이에 전용 암호화 지하 터널을 만든다고 상상해 보세요. 터널 안에서 걸어다니면, 바깥의 사람들은 여러분이 무엇을 들고 있는지, 어디로 가는지 볼 수 없습니다.

\begin{itemize}
  \item \textbf{일반 인터넷 사용}: 스마트폰 공용 네트워크 대상 웹사이트. 중간에 거치는 라우터, 스위치 등 어느 단계에서든 누군가가 여러분이 보낸 데이터를 "훔쳐볼" 수 있습니다. 특히 카페, 공항 같은 공용 Wi-Fi에서는 거칠게 말해 "벌거벗고 다니는" 것과 다름없습니다.
  \item \textbf{VPN 사용}: 스마트폰 VPN 암호화 터널 VPN 서버 대상 웹사이트. 데이터가 암호화되어 캡슐화되므로, "엿보는 사람"은 암호화된 문자만 볼 수 있습니다. 또한 대상 웹사이트가 보는 것은 VPN 서버의 IP 주소이며, 여러분의 실제 IP가 아닙니다.
\end{itemize}

VPN이 실제로 무엇을 할 수 있나요? 간단히 말하면: 첫째,\textbf{공용 Wi-Fi에서의 프라이버시 보호}; 둘째,\textbf{지역 제한 우회}(예: 특정 국가에서만 출시된 콘텐츠 시청); 셋째,\textbf{인터넷 서비스 제공업체(ISP)가 브라우징 기록을 기록하는 것 방지}. 다만 한 가지 알려드립니다:\textbf{무료 VPN은 사용하지 마세요!}무료의 이면에는 종종 여러분의 데이터가 그들의 “상품”이라는 의미입니다. 지인이 어떤 무료 VPN을 사용했는데, 스마트폰에 이유 없이 광고가 팝업되기 시작했고, 알고 보니 그 VPN이 브라우징 데이터를 모아 광고 회사에 팔아넘긴 것이었습니다. 따라서 VPN을 선택하려면 유료의 대기업 브랜드(예: NordVPN, ExpressVPN)를 선택하거나, 차라리 사용하지 않는 것이 좋습니다.

또한, VPN이 만능은 아닙니다. “전송 중”인 데이터만 보호할 수 있으며, 대상 웹사이트 자체가 안전하지 않거나, 스마트폰에 이미 악성 소프트웨어가 설치되어 있다면, VPN도 도움이 되지 못합니다.

\booksection{둘, 공용 Wi-Fi의 “달콤한 함정” — 상상보다 위험하다}

그냥 무료 Wi-Fi에 연결하는 거라고 생각할 수 있죠. 얼마나 위험할 수 있겠냐고요? 실제 사례를 알려드리겠습니다: 보안 연구원들이 실험을 한 적이 있는데, 공항에 “Starbucks\_Free”라는 Wi-Fi 핫스팟을 설치했습니다. 1시간도 안 되어 수백 명이 연결했습니다. 사실 이 핫스팟은 연구원들이 자체적으로 구축한 것이었고, 이 핫스팟을 통과하는\textbf{암호화되지 않은 트래픽}— 예를 들어 방문한 HTTP 웹페이지 내용, TLS가 활성화되지 않은 앱 통신 — 을 실제로 가로채고 확인할 수 있었습니다.하지만 핵심 전제를 짚어야 합니다: 오늘날 일상적으로 사용하는 위챗, 인터넷 뱅킹, Alipay는 대부분 이미 HTTPS 암호화를 사용하고 있습니다. 솔직히 말해, 공격자가 데이터 패킷을 가로채도 보이는 것은 암호화된 문자뿐이며,\textbf{채팅 내용이나 결제 비밀번호를 직접 읽을 수는 없습니다}. 하지만 여전히 어떤 웹사이트를 방문했는지(도메인 이름)는 알 수 있습니다. 일부 오래된 앱이나 웹페이지가 여전히 HTTP 평문으로 전송한다면, 그 내용이 완전히 노출됩니다. 이러한 공격을 “중간자 공격(Man-in-the-Middle)”이라고 하며, 원리는 공격자가 여러분과 진짜 Wi-Fi 라우터 사이에 “끼어드는” 것입니다. 더 무서운 것은, 많은 사람들의 스마트폰이 “알려진 Wi-Fi에 자동 연결”로 설정되어 있어, 공격자가 자신의 핫스팟 이름을 이전에 연결했던 Wi-Fi 이름(예: “ChinaNet” 또는 “CMCC”)으로 설정하면, 스마트폰이 자동으로 연결됩니다.그러면 어떻게 방어하나요? 몇 가지 실용적인 조언:\textbf{HTTPS 식별}: 주소 표시줄이 https://로 시작하며, 추가된 “s”는 Secure를 의미합니다. 비밀번호를 입력해야 하는 웹사이트를 방문할 때, http://(s가 없음)이 보이면 즉시 입력을 멈추세요! 요즘 주요 브라우저는 HTTPS가 아닌 웹사이트 앞에 “안전하지 않음” 경고를 표시하며, 이 표시를 보면 경계해야 합니다.\textbf{“자동 연결” 끄기}: iPhone이든 안드로이드든, Wi-Fi 설정에서 “자동 가입” 또는 “자동 연결”을 끄세요. 매번 수동으로 한 번 클릭하더라도, 스마트폰이 스스로 판단하지 않도록 하세요.\textbf{민감한 작업에는 스마트폰 데이터 사용}: 계좌 이체, 인터넷 뱅킹 로그인, 주민등록번호 입력... 이런 작업은 성실히 자신의 4G/5G 데이터를 사용하고, 약간의 데이터 요금을 아끼려 하지 마세요.\textbf{자주 사용하지 않는 Wi-Fi 잊어버리기}: 출장, 여행에서 돌아온 후, 이전에 임시로 연결했던 호텔 Wi-Fi, 공항 Wi-Fi를 모두 “잊어버리세요”. 이렇게 하면 다음에 스마트폰이 같은 이름의 가짜 핫스팟에 자동으로 연결되지 않습니다.셋, HTTPS는 어떻게 여러분을 보호하는가? — SSL/TLS 핸드셰이크앞서 HTTPS를 언급했는데, 실제로 어떻게 작동할까요? HTTPS 웹사이트를 방문할 때마다, 브라우저와 서버 사이에 “SSL/TLS 핸드셰이크”가 한 번 수행되며, 전체 과정은 대략 다음과 같습니다:브라우저: “이봐, 서버, 나는 안전하게 통신하고 싶어, 이건 내가 지원하는 암호화 능력 목록이야(지원하는 암호화 알고리즘 등).”서버 응답: “알겠어, 나는 당신의 목록에서 이 암호화 방식을 선택할게, 이건 내 디지털 인증서이며, 내 공개 키가 포함되어 있고, 모 CA 기관의 서명 인증을 받았어.”브라우저가 인증서의 진위 검증: 신뢰할 수 있는 CA에서 발급했는지? 만료되지 않았는지? 도메인 이름이 일치하는지?검증이 통과되면, 양측은 공개 키로 암호화하여 “세션 키”를 협상하며, 이후의 통신은 모두 이 키로 암호화됩니다.은행에 가서 업무를 보는 것과 같습니다: 먼저 신분증을 제시하면(인증서), 은행원이 신원을 확인한 후 둘이 암호를 정하고, 이후의 대화는 둘만 이해할 수 있습니다. 옆에서 누가 엿들어도 암호화된 문자만 들을 수 있습니다.하지만 주의: HTTPS는 “전송 과정”이 암호화되었음을 보장할 뿐, 대상 웹사이트가 좋은 사이트임을 보장하지 않습니다. 피싱 웹사이트도 HTTPS를 사용할 수 있습니다 — 그 작은 자물쇠 아이콘은 단지 여러분이 “어떤 웹사이트”에 연결되어 있음을 의미할 뿐, 그 웹사이트가 여러분이 가고 싶었던 곳이라는 의미는 아닙니다. 따라서 여전히 웹사이트 주소가 정확한지 직접 확인해야 합니다.넷, DNS 보안: “인터넷 전화번호부”가 안전한가?DNS(Domain Name System)는 인터넷의 전화번호부와 같습니다: “baidu.com”을 입력하면, DNS가 해당 IP 주소(예: 110.242.68.66)를 조회해 주고, 그래야 브라우저가 연결할 수 있습니다. 문제는 전통적인 DNS 조회가\textbf{평문}이고,\textbf{암호화되지 않았다}는 것입니다! 즉, ISP, Wi-Fi 제공자 또는 중간 노드가 여러분이 어떤 웹사이트를 방문했는지 볼 수 있습니다.더 끔찍한 것은 “DNS 하이재킹”입니다: 공격자가 DNS 단계에서 조작하여 “baidu.com”을 가짜 피싱 웹사이트로 연결합니다. 분명 올바른 웹사이트 주소를 입력했는데, 똑같이 생겼지만 실제로는 가짜인 로그인 페이지로 이동하여, 계정과 비밀번호를 입력하는 순간 끝납니다.다행히 이제 암호화 DNS 방안이 있습니다:\textbf{DNS over HTTPS(DoH)}: DNS 조회를 HTTPS 트래픽 안에 숨겨, 정상적인 웹페이지 브라우징 트래픽과 섞어 DNS 조회인지 구별할 수 없게 합니다.\textbf{DNS over TLS(DoT)}: TLS 암호화 프로토콜로 DNS 조회를 보호하며, 포트가 독립적이라, 관리자가 DoT를 사용하는 것은 알 수 있지만 내용은 볼 수 없습니다.개인적으로 Cloudflare의 1.1.1.1 또는 Google의 8.8.8.8을 DNS 서버로 사용하는 것을 권장하며, 둘 다 암호화 DNS를 지원합니다. 안드로이드의 “개인 DNS” 설정에서 “one.one.one.one” 또는 “dns.google”을 직접 입력하면 됩니다. 매우 간단합니다. iPhone은 약간 번거로운데, 프로파일을 설치하거나 앱을 통해 설정해야 하지만, 그럴 가치는 있습니다.다섯, “시크릿 모드”는 정말 투명인간이 될 수 있는가?먼저 결론을 말합니다:\textbf{시크릿 모드는 집의 커튼이지 투명 망토가 아닙니다.}하는 일: 기기에 흔적을 남기지 않음 — 시크릿 창을 닫으면, 브라우저는 브라우징 기록, 쿠키, 양식 작성 기록, 검색 기록을 저장하지 않습니다. 남의 컴퓨터로 뭔가를 검색하거나, 선물을 사서 가족이 브라우징 기록으로 발견하지 않기를 바라는 경우, 시크릿 모드를 사용하는 것이 적합합니다.하지 못하는 일:방문한 웹사이트는\textbf{여전히 여러분이 방문한 것을 알 수 있으며}, IP 주소, 기기 정보, 브라우징 행동 등 웹사이트 측에서 명확히 파악합니다.인터넷 서비스 제공업체는\textbf{여전히 어떤 웹사이트를 방문했는지 알 수 있으며}, (HTTPS를 사용한 경우) 구체적인 내용은 볼 수 없습니다.회사/학교 네트워크 관리자는\textbf{동일하게 모니터링}할 수 있습니다.다운로드한 파일과 추가한 북마크는\textbf{유지되며}, 창을 닫아도 사라지지 않습니다.따라서 기억하세요: 시크릿 모드는 “로컬 지우개”이지 “네트워크 투명 망토”가 아닙니다. 만약 인터넷 흔적을 진짜로 숨기고 싶다면, VPN+Tor 브라우저와 같은 조합이 필요하지만, 일상적인 사용에는 이 정도만 알아도 충분합니다.여섯, 브라우저 프라이버시: 브라우저가 여러분을 팔지 않도록의식하지 못했을 수 있지만, 브라우저 자체가 거대한 정보 수집기입니다. 쿠키, 브라우징 기록, 캐시, 자동 완성, 위치 권한... 모든 것이 프라이버시를 유출할 수 있습니다. 몇 가지 조언:\textbf{주기적으로 브라우징 데이터 정리}: 가끔 설정에서 쿠키, 캐시, 기록을 정리하세요. “추적 방지”를 위한 것이 아니라 — 쿠키 정리만으로는 이미 효과적이지 않습니다 — 스마트폰이 점점 느려지는 것을 방지하기 위함입니다. 브라우저 캐시가 수 GB의 공간을 차지할 수 있습니다.\textbf{“추적 거부” 환상 버리기}: 브라우저의 “추적 거부 요청 보내기” 옵션은 본질적으로 웹사이트에 “추적하지 말아 달라”고 “요청”하는 것입니다. 하지만 웹사이트가 따르느냐 마느냐는 완전히 다른 문제이며, 대부분의 웹사이트는 직접 무시합니다. 기대하지 마세요.\textbf{프라이버시 보호 브라우저 활용}: Firefox Focus, DuckDuckGo 브라우저, Brave는 기본적으로 추적기와 광고를 차단하므로, Chrome과 Safari의 기본 설정보다 훨씬 깨끗합니다.\textbf{검색 엔진도 중요}: Baidu/Google 검색 대신 DuckDuckGo를 사용하세요. 검색 기록을 저장하지 않고, 사용자 프로필을 구축하지 않습니다. 솔직히 검색 결과는 Google만큼 정확하지는 않지만, 일상적인 자료 조회에는 충분합니다.\textbf{앱 권한 확인}: 스마트폰 설정에서 브라우저에 어떤 권한이 부여되어 있는지 확인하세요. 위치, 마이크, 주소록에 정말로 접근할 필요가 있나요? 아마 필요 없을 것입니다.일곱, 실용적인 보안 체크리스트: 지금부터 자신 보호이야기가 많았는데, 마지막으로 간단하고 강력한 보안 체크리스트를 드립니다. 30분 정도만 투자하면, 스마트폰 보안 수준이 질적으로 향상될 것입니다:\textbf{Wi-Fi 자동 연결 끄기}: iPhone은 설정-무선 LAN에서, 안드로이드는 Wi-Fi 설정에서.\textbf{암호화 DNS로 변경}: 안드로이드는 개인 DNS를 “dns.google”로 직접 설정하고, iPhone은 앱 또는 프로파일을 통해 1.1.1.1의 DoH를 설정합니다.\textbf{2단계 인증(2FA) 활성화}: 위챗, Alipay, Apple ID, Google 계정 등 2FA를 켤 수 있는 것은 모두 켜세요. SMS 인증 코드는 이미 최저 수준의 2FA이며, TOTP(예: Google Authenticator, Microsoft Authenticator)를 사용할 수 있다면 더 좋습니다.\textbf{앱 권한 한 번 확인}: 각 앱의 권한을 한 번 검토하세요, 특히 위치, 주소록, SMS, 마이크, 카메라의 5가지 항목. 손전등 앱이 주소록을 요구하는 경우가 있다는 것을 알게 될 것입니다 — 즉시 삭제하세요.\textbf{주기적으로 브라우저 데이터 정리}: 캘린더 알림을 설정하여, 매월 한 번 캐시, 쿠키를 정리하세요.\textbf{공용 충전 스테이션의 데이터 케이블 사용 금지}: USB 포트는 충전과 데이터 전송이 모두 가능하며, 악성 충전 스테이션(“Juice Jacking”이라고 함)은 충전 중에 데이터를 훔칠 수 있습니다. 외출 시 자신의 충전기와 케이블을 휴대하세요.지금 하기: 스마트폰에 “보안 건강 검진” 실시\textbf{DNS 업그레이드}: 안드로이드: 설정 > 연결 및 공유 > 개인 DNS > “one.one.one.one” 또는 “dns.google” 입력. iPhone: Safari에서 “1.1.1.1 DNS profile” 검색하여 프로파일 설치, 또는 Cloudflare의 1.1.1.1 앱 다운로드.\textbf{Wi-Fi 감사}: Wi-Fi 목록을 열어, 모르거나 자주 사용하지 않거나 이름이 이상한 핫스팟을 모두 “잊어버리세요”.\textbf{권한 대청소}: 설정 > 프라이버시/권한 관리로 들어가, 앱별로 권한을 검토. 권한 요구가 불합리한 앱(예: 계산기가 주소록 요구)은 표시하고, 삭제하거나 대안을 고려.\textbf{계정 보안 확인}: Apple ID 또는 Google 계정의 보안 페이지를 열어, 알 수 없는 기기 로그인이 있는지 확인하고, 2FA가 활성화되었는지 확인하며, 복구 연락처가 신뢰할 수 있는지 확인.
